\documentclass[10pt,a4paper]{ctexart}

\usepackage[margin=1.65cm]{geometry}
\usepackage{amsmath,amssymb}
\usepackage{booktabs,longtable,tabularx,array,multicol}
\usepackage{enumitem}
\usepackage{xcolor}
\usepackage{hyperref}
\usepackage{fancyhdr}
\usepackage{makeidx}

\definecolor{Ink}{HTML}{1F2937}
\definecolor{Muted}{HTML}{5B6673}
\definecolor{Teal}{HTML}{007C7C}
\definecolor{Indigo}{HTML}{243B6B}
\definecolor{Amber}{HTML}{B86B00}
\definecolor{Soft}{HTML}{F4F7F6}

\hypersetup{
  colorlinks=true,
  linkcolor=Indigo,
  urlcolor=Teal,
  citecolor=Amber
}

\pagestyle{fancy}
\fancyhf{}
\lhead{Open-book Exam Guide}
\rhead{开卷考中英资料}
\cfoot{\thepage}
\setlength{\headheight}{14pt}

\setlength{\parindent}{0pt}
\setlength{\parskip}{3pt}
\setlist[itemize]{topsep=2pt,itemsep=1pt,leftmargin=1.25em}
\setlist[enumerate]{topsep=2pt,itemsep=1pt,leftmargin=1.45em}
\renewcommand{\arraystretch}{1.16}
\newcolumntype{L}[1]{>{\raggedright\arraybackslash}p{#1}}
\newcommand{\term}[2]{\textbf{\textcolor{Teal}{#1}} \textcolor{Muted}{(#2)}}
\newcommand{\source}[1]{\textcolor{Muted}{\footnotesize Source: #1}}
\newcommand{\exam}[1]{\textcolor{Amber}{\textbf{Exam use:} #1}}
\newcommand{\cn}[1]{\textbf{中文：}#1}
\newcommand{\en}[1]{\textbf{English:} #1}
\newcommand{\kw}[1]{\textbf{\textcolor{Indigo}{#1}}}
\makeindex

\title{\textbf{Operations Research and Information System Management}\\
\large 第一次开卷考资料 - 中英双语速查版}
\author{Based on OR Class 1-13 slides}
\date{Prepared for open-book review}

\begin{document}
\maketitle
\tableofcontents
\newpage

\section*{使用方法 / How to Use This Guide}
\addcontentsline{toc}{section}{使用方法 / How to Use This Guide}

\begin{itemize}
  \item \cn{开卷考最重要的是定位速度。先看第 \ref{sec:fast-index} 节的索引表，用题目关键词定位到公式、步骤和对应课件。}
  \item \en{For an open-book exam, retrieval speed matters most. Start from Section \ref{sec:fast-index}: use keywords to locate formulas, procedures, and slide sources.}
  \item \cn{计算题优先看第 \ref{sec:templates} 节的题型模板；概念题优先看各章的 ``must know'' 和 ``pitfall''。}
  \item \en{For calculation questions, use the templates in Section \ref{sec:templates}. For concept questions, use the ``must know'' and ``pitfall'' notes in each topic section.}
  \item \cn{末尾有英文术语索引。前置索引更适合考试，末尾索引更适合按英文关键词查。}
  \item \en{The back index is sorted by English terms. The front index is more useful during the exam.}
\end{itemize}

\section{开卷检索索引 / Fast Retrieval Index}
\label{sec:fast-index}

\subsection{题目关键词索引 / Question Keyword Index}

{\footnotesize
\setlength{\tabcolsep}{3pt}
\begin{longtable}{@{}L{0.17\linewidth}L{0.28\linewidth}L{0.29\linewidth}L{0.13\linewidth}@{}}
\toprule
\textbf{看到关键词 / If you see} & \textbf{马上找 / Go to} & \textbf{答题骨架 / Answer skeleton} & \textbf{Slides} \\
\midrule
forecast, demand, trend, alpha, smoothing, moving average & 预测 / Forecasting, Section \ref{sec:forecasting} & 选方法 $\rightarrow$ 算预测值 $\rightarrow$ 算误差 $\rightarrow$ 比较 MAD/MSE/MAPE & Class 4, 5 \\
MAD, MSE, MAPE, tracking signal, control chart & 预测误差 / Accuracy, Section \ref{subsec:accuracy} & $e=A-F$ $\rightarrow$ 绝对值/平方/百分比 $\rightarrow$ 判断 bias & Class 5 \\
regression, least squares, fitted line, trend projection & 趋势与回归 / Trend, Section \ref{subsec:trend} & 算 $a,b$ $\rightarrow$ 写 $\hat{y}=a+bt$ 或 $\hat{y}=a+bx$ $\rightarrow$ 代入预测 & Class 4, 5 \\
maximin, maximax, Laplace, regret & 不确定性决策 / Uncertainty, Section \ref{subsec:uncertainty} & 收益表 $\rightarrow$ 每个方案求对应准则 $\rightarrow$ 选最佳 & Class 9 \\
probability, risk, EMV & 风险决策 / EMV, Section \ref{subsec:emv} & 每方案 $EMV=\sum p_s payoff$ $\rightarrow$ 选最大 & Class 9 \\
decision tree, backward induction & 决策树 / Decision tree, Section \ref{subsec:tree} & 从左往右读，从右往左算；机会节点算期望，决策节点取最大 & Class 9 \\
EVPI, perfect information & 完全信息价值 / EVPI, Section \ref{subsec:evpi} & $A=\sum p_s \max payoff$，$B=\max EMV$，$EVPI=A-B$ & Class 9 \\
capacity, utilization, efficiency, design capacity, effective capacity & 产能度量 / Capacity measures, Section \ref{subsec:cap-measures} & 分清三个 capacity $\rightarrow$ efficiency 用 effective，utilization 用 design & Class 10 \\
capacity cushion, lead/follow/track & 产能战略 / Capacity strategy, Section \ref{subsec:cap-strategy} & 不确定性越高，cushion 越大；三种扩张策略比较 & Class 10 \\
bottleneck, constraint & 瓶颈与约束 / Bottleneck, Section \ref{subsec:bottleneck} & 系统产出由最慢环节限制；先改善 bottleneck & Class 11, 13 \\
BEP, fixed cost, variable cost, make or buy & 成本-产量分析 / Cost-volume, Section \ref{subsec:cvp} & 写 $TC,TR,Profit$ $\rightarrow$ 求交点或目标利润产量 & Class 11 \\
standardization, delayed differentiation, modular design & 产品设计策略 / Design strategy, Section \ref{subsec:standardization} & 标准化提高效率；延迟差异化推迟定制；模块化兼顾标准化和定制 & Class 6 \\
reliability, robust design, Taguchi & 可靠性与稳健设计 / Reliability, Section \ref{subsec:reliability} & reliability 是规定条件下正常工作；robust 是不同环境下保持稳定 & Class 6 \\
FMEA, FTA, VA, QFD, house of quality, Kano & 质量设计工具 / Quality tools, Section \ref{subsec:quality-tools} & 工具名称 $\rightarrow$ 目的 $\rightarrow$ 表格/矩阵/分类逻辑 & Class 7 \\
service blueprint, high-contact, low-contact & 服务设计 / Service design, Section \ref{subsec:service-design} & 服务概念/服务包/服务规格；蓝图步骤；高低接触比较 & Class 8 \\
job shop, batch, repetitive, continuous, project & 流程选择 / Process types, Section \ref{subsec:process-types} & 产量低高、品种多少、柔性高低三维定位 & Class 12 \\
product layout, process layout, fixed-position, cellular & 设施布局 / Layouts, Section \ref{subsec:layouts} & product 顺序流；process 功能分组；fixed 产品不动；cellular 按族群 & Class 12, 13 \\
line balancing, cycle time, workstation, idle time & 生产线平衡 / Line balancing, Section \ref{subsec:line-balancing} & 算 $CT$ $\rightarrow$ 算 $N_{\min}$ $\rightarrow$ 按前置关系分配 $\rightarrow$ idle/efficiency & Class 13 \\
transportation cost, flow, distance, process layout assignment & 流程布局成本 / Load-distance, Section \ref{subsec:load-distance} & 大 flow 配短 distance；总成本 $\sum F_{ij}D_{ij}C_{ij}$ & Class 13 \\
\bottomrule
\end{longtable}
}

\subsection{公式索引 / Formula Index}

\begin{longtable}{@{}L{0.25\linewidth}L{0.36\linewidth}L{0.30\linewidth}@{}}
\toprule
\textbf{公式 / Formula} & \textbf{表达 / Expression} & \textbf{什么时候用 / When to use} \\
\midrule
Productivity\index{Productivity} & $\text{Productivity}=\frac{\text{Output}}{\text{Input}}$ & 生产率、资源利用效率。 \\
Productivity growth & $\frac{P_{current}-P_{previous}}{P_{previous}}$ & 比较两期生产率增长。 \\
Naive forecast\index{Naive forecast} & $F_t=A_{t-1}$ & 最简单预测，下一期等于上一期实际值。 \\
Moving average\index{Moving average} & $F_t=\frac{A_{t-1}+\cdots+A_{t-n}}{n}$ & 稳定序列，平滑随机波动。 \\
Weighted moving average\index{Weighted moving average} & $F_t=\sum w_i A_{t-i}$, $\sum w_i=1$ & 近期数据更重要。 \\
Exponential smoothing\index{Exponential smoothing} & $F_t=F_{t-1}+\alpha(A_{t-1}-F_{t-1})$ & 只需上一期 forecast 和 actual。 \\
Linear trend\index{Linear trend} & $F_t=a+bt$ & 有线性趋势时预测。 \\
Least-squares slope & $b=\frac{n\sum ty-\sum t\sum y}{n\sum t^2-(\sum t)^2}$ & 估计趋势线斜率。 \\
Least-squares intercept & $a=\frac{\sum y-b\sum t}{n}$ & 估计趋势线截距。 \\
MAD\index{MAD} & $MAD=\frac{\sum |A_t-F_t|}{n}$ & 平均绝对误差。 \\
MSE\index{MSE} & $MSE=\frac{\sum (A_t-F_t)^2}{n-1}$ & 大误差惩罚更重。 \\
MAPE\index{MAPE} & $MAPE=\frac{100}{n}\sum\left|\frac{A_t-F_t}{A_t}\right|$ & 百分比误差，可跨规模比较。 \\
Tracking signal\index{Tracking signal} & $\frac{\sum(A_t-F_t)}{MAD}$ & 判断持续偏高或偏低。 \\
EMV\index{EMV} & $EMV_i=\sum_s p_s payoff_{i,s}$ & 风险下选最大期望收益。 \\
EVPI\index{EVPI} & $EVPI=\sum_s p_s\max_i payoff_{i,s}-\max_i EMV_i$ & 完全信息的最高价值。 \\
Efficiency\index{Efficiency} & $\frac{\text{Actual Output}}{\text{Effective Capacity}}$ & 衡量有效产能使用程度。 \\
Utilization\index{Utilization} & $\frac{\text{Actual Output}}{\text{Design Capacity}}$ & 衡量设计产能使用程度。 \\
Capacity cushion\index{Capacity cushion} & $100\%-Utilization$ & 应对需求不确定性的额外产能。 \\
Processing requirement & $NR=\frac{\sum p_iD_i}{T}$ & 多产品加工资源需求。 \\
Total cost\index{Cost-volume analysis} & $TC=FC+vcQ$ & 成本-产量分析。 \\
Total revenue & $TR=RQ$ & 收入-产量分析。 \\
Profit & $P=Q(R-vc)-FC$ & 目标利润或盈亏分析。 \\
BEP\index{Break-even point} & $Q_{BEP}=\frac{FC}{R-vc}$ & 盈亏平衡点。 \\
Cycle time\index{Cycle time} & $CT=\frac{OT}{D}$ & 生产线平衡，目标产出给定。 \\
Minimum stations & $N_{\min}=\left\lceil \frac{\sum t_i}{CT}\right\rceil$ & 最少工作站数量。 \\
Idle time & $\frac{\text{Idle time per cycle}}{N\times CT}$ & 生产线闲置比例。 \\
Line efficiency\index{Line balancing} & $1-\text{Percent idle time}$ & 生产线效率。 \\
Load-distance cost\index{Load-distance cost} & $\sum_{i<j}F_{ij}D_{ij}C_{ij}$ & 流程布局运输成本。 \\
\bottomrule
\end{longtable}

\section{运营管理基础、战略与生产率 / OM Basics, Strategy, and Productivity}
\label{sec:om}

\subsection{核心定义 / Core Definitions}

\term{Operations Management}{运营管理}\index{Operations management}: management of the transformation process that turns inputs into outputs. \\
\cn{运营管理关注如何把输入转化为产品、服务或信息，并创造价值。}\\
\en{Operations management focuses on converting inputs into goods, services, or information that create value.}\\
\source{OR Class 2, slides 3-10}

\textbf{Input - Transformation - Output}
\begin{itemize}
  \item Inputs: labor, equipment, facilities, raw materials, time, knowledge, skills, customer relationships, reputation.
  \item Outputs: products, services, information.
  \item Transformation can be physical, locational, exchange, storage, physiological, or informational.
\end{itemize}

\subsection{系统设计 vs. 系统运行 / System Design vs. System Operation}

\begin{longtable}{@{}L{0.46\linewidth}L{0.46\linewidth}@{}}
\toprule
\textbf{System Design Decisions / 系统设计决策} & \textbf{System Operation Decisions / 系统运行决策} \\
\midrule
Capacity, location, layout, product/service planning, acquisition and placement of equipment. & Personnel management, inventory control, scheduling, project management, quality assurance. \\
Strategic, long-term, resource commitment. & Tactical or operational, day-to-day control. \\
\bottomrule
\end{longtable}

\exam{如果题目问 product design 和 process selection 属于哪类决策，答案是 system design decisions。}

\subsection{模型、权衡与系统观 / Models, Trade-offs, and Systems Approach}
\term{Model}{模型}\index{Model}: abstraction and simplification of reality. Models help organize information, quantify problems, and answer what-if questions.\\
\term{Trade-off}{权衡}\index{Trade-off}: giving up one thing in exchange for another, such as cost vs. quality, speed vs. flexibility.\\
\term{Systems approach}{系统观}\index{Systems approach}: emphasizes relationships among subsystems; the whole is greater than the sum of parts.\\
\exam{模型题常考 ``models do not guarantee the best decision''。模型有用，但会遗漏定性因素，也可能被误用。}

\subsection{竞争力、订单资格和赢单要素 / Competitiveness, Qualifiers, and Winners}

\term{Competitiveness}{竞争力}\index{Competitiveness}: how effectively an organization meets customer wants and needs relative to competitors.\\
\term{Order qualifiers}{订单资格要素}\index{Order qualifiers}: minimum standards that allow an option to be considered.\\
\term{Order winners}{赢单要素}\index{Order winners}: characteristics that make customers choose one option over competitors.\\
\source{OR Class 3, slides 5, 31-32}

\textbf{Strategy formulation / 战略形成}
\begin{enumerate}
  \item Define the primary task.
  \item Assess distinctive competencies.
  \item Identify order qualifiers and order winners.
  \item Position the firm.
\end{enumerate}

\subsection{生产率 / Productivity}
\index{Productivity}
\[
\text{Productivity}=\frac{\text{Outputs}}{\text{Inputs}}, \qquad
\text{Growth}=\frac{P_{current}-P_{previous}}{P_{previous}}
\]

\begin{itemize}
  \item \term{Partial productivity}{单要素生产率}: output divided by one input, e.g. units per labor hour.
  \item \term{Multifactor productivity}{多要素生产率}: output divided by several inputs.
  \item \term{Total productivity}{总生产率}: output divided by all inputs.
\end{itemize}
\exam{题目给 labor/material/overhead 时，labor productivity 只用 labor；total productivity 用所有 input 的总和。}

\section{预测 / Forecasting}
\label{sec:forecasting}

\subsection{预测基本逻辑 / Forecasting Logic}
\term{Forecast}{预测}\index{Forecasting}: a statement about the future value of a variable of interest, usually demand.\\
\source{Class 4, slides 6-11}

\textbf{Good forecasts / 好预测的特征}
\begin{multicols}{2}
\begin{itemize}
  \item Timely / 及时
  \item Accurate / 准确
  \item Reliable / 可靠
  \item Meaningful units / 单位有意义
  \item Written / 书面化
  \item Easy to understand / 易理解
  \item Cost-effective / 成本有效
\end{itemize}
\end{multicols}

\textbf{Forecasting process / 预测步骤}
\begin{enumerate}
  \item Determine purpose.
  \item Establish time horizon.
  \item Select technique.
  \item Gather and analyze data.
  \item Make forecast.
  \item Monitor forecast.
\end{enumerate}

\subsection{定性预测 / Qualitative Forecasting}

\begin{longtable}{@{}L{0.22\linewidth}L{0.36\linewidth}L{0.34\linewidth}@{}}
\toprule
\textbf{Method} & \textbf{中文解释} & \textbf{Exam note} \\
\midrule
Executive opinions & 高层小组判断。 & 快，但可能受权力和偏见影响。 \\
Sales force opinions & 销售人员接触客户。 & 信息直接，但可能有激励偏差。 \\
Consumer survey & 问卷调查顾客意向。 & 新产品有用，成本高。 \\
Delphi method\index{Delphi method} & 多轮匿名专家问卷。 & 减少从众和权威压力。 \\
Outside opinions & 第三方预测。 & 省力，但依赖外部机构质量。 \\
\bottomrule
\end{longtable}

\subsection{时间序列方法 / Time-Series Methods}
\index{Time series}

\begin{align*}
\text{Naive:}\quad & F_t=A_{t-1}\\
\text{Moving Average:}\quad & F_t=\frac{A_{t-1}+\cdots +A_{t-n}}{n}\\
\text{Weighted Moving Average:}\quad & F_t=\sum_{i=1}^{n}w_iA_{t-i}, \quad \sum w_i=1\\
\text{Exponential Smoothing:}\quad & F_t=F_{t-1}+\alpha(A_{t-1}-F_{t-1})
\end{align*}

\textbf{Alpha interpretation / $\alpha$ 的含义}\index{Alpha}
\begin{itemize}
  \item Small $\alpha$: stable, smooth, slow response, larger lag.
  \item Large $\alpha$: responsive, tracks actual series more closely, less smoothing.
  \item Slides note: common rule of thumb $\alpha <0.5$; $\alpha=0.2$ or $0.3$ often works.
\end{itemize}
\source{Class 4, slides 44-49}

\subsection{趋势与回归 / Trend and Regression}
\label{subsec:trend}
\index{Regression}
\[
F_t=a+bt
\]
\[
b=\frac{n\sum ty-\sum t\sum y}{n\sum t^2-(\sum t)^2},
\qquad
a=\frac{\sum y-b\sum t}{n}
\]
\[
SSE=\sum_{i=1}^{n}(y_i-\hat{y}_i)^2
\]

\textcolor{Amber}{\textbf{Exam use:}} 如果题目给 period/week/month 与 sales/demand，并问 least squares trend line：
\begin{enumerate}
  \item 建表求 $\sum t,\sum y,\sum t^2,\sum ty$。
  \item 计算 $b,a$。
  \item 把未来期数代入 $F_t=a+bt$。
\end{enumerate}

\subsection{预测误差 / Forecast Accuracy}
\label{subsec:accuracy}
\index{Forecast accuracy}

\[
e_t=A_t-F_t
\]
\[
MAD=\frac{\sum |A_t-F_t|}{n}
\]
\[
MSE=\frac{\sum (A_t-F_t)^2}{n-1}
\]
\[
MAPE=\frac{100}{n}\sum\left|\frac{A_t-F_t}{A_t}\right|
\]
\[
\text{Tracking Signal}=\frac{\sum(A_t-F_t)}{MAD}
\]

\begin{longtable}{@{}L{0.20\linewidth}L{0.38\linewidth}L{0.34\linewidth}@{}}
\toprule
\textbf{Measure} & \textbf{中文} & \textbf{English} \\
\midrule
MAD & 平均绝对误差，容易解释。 & Average absolute error; linear penalty. \\
MSE & 平方误差均值，对大误差惩罚更重。 & Penalizes large errors more strongly. \\
MAPE & 百分比误差，适合跨规模比较。 & Percentage error; scale-free comparison. \\
Tracking signal & 累计误差除以 MAD，判断系统偏差。 & Detects persistent bias. \\
\bottomrule
\end{longtable}

\exam{题目说 ``MAD is the square root of MSE'' 是 false。MAD 是 mean absolute deviation，不是 MSE 的平方根。}

\section{产品与服务设计 / Product and Service Design}
\label{sec:design}

\subsection{设计目标 / Design Objectives}
\index{Product design}
\cn{产品和服务设计的核心是把顾客需求转化为可运营、可制造、可交付的规格。}\\
\en{The core of product and service design is translating customer needs into operational, manufacturable, and deliverable specifications.}\\
\source{Class 6, slides 8-16}

\textbf{Main issues / 主要问题}
\begin{multicols}{2}
\begin{itemize}
  \item Economic / 经济
  \item Competitive / 竞争
  \item Legal / 法律
  \item Ethical / 伦理
  \item Environmental / 环境
  \item Operational capabilities / 运营能力
\end{itemize}
\end{multicols}

\subsection{标准化、延迟差异化、模块化 / Standardization, Delayed Differentiation, Modularity}
\label{subsec:standardization}
\index{Standardization}
\index{Delayed differentiation}
\index{Modular design}

\begin{longtable}{@{}L{0.24\linewidth}L{0.36\linewidth}L{0.32\linewidth}@{}}
\toprule
\textbf{Concept} & \textbf{中文解释} & \textbf{Exam note} \\
\midrule
Standardization & 减少产品、服务或流程的变化。 & 优点是低成本、一致性、长生产批量；缺点是缺少多样性，设计变化阻力大。 \\
Component commonality & 多个产品使用相同零件。 & 降低库存、采购、培训复杂度。 \\
Delayed differentiation & 先生产但不最终完成，等顾客偏好明确后再定制。 & 是 postponement tactic。 \\
Modular design & 把组件分成可替换模块。 & 便于诊断、维修、替换、定制；可能增加整模块替换成本。 \\
Mass customization & 标准化基础上的定制。 & 用标准化效率支持个性化。 \\
\bottomrule
\end{longtable}

\exam{题目问 postponement tactic 的例子，选 delayed differentiation。}

\subsection{可靠性与稳健设计 / Reliability and Robust Design}
\label{subsec:reliability}
\index{Reliability}
\index{Robust design}

\term{Reliability}{可靠性}: ability of a product, part, or system to perform its intended function under prescribed conditions.\\
\term{Failure}{失效}: product, part, or system does not perform as intended.\\
\term{Robust design}{稳健设计}: product or service functions over a broad range of conditions.\\
\source{Class 6, slides 44-49}

\textbf{Improve reliability / 提升可靠性}
\begin{multicols}{2}
\begin{itemize}
  \item Component design
  \item Production/assembly techniques
  \item Testing
  \item Redundancy/backup
  \item Preventive maintenance
  \item User education
  \item System design
\end{itemize}
\end{multicols}

\exam{题目问 one way to increase reliability，improve component design 或 add backup components 都符合。}

\subsection{产品开发流程与并行工程 / Development Process and Concurrent Engineering}
\index{Concurrent engineering}
\index{Design for manufacturing}

\textbf{Product development process / 产品开发流程}
\begin{enumerate}
  \item Idea generation / 创意产生
  \item Feasibility analysis / 可行性分析
  \item Product specifications / 产品规格
  \item Process specifications / 流程规格
  \item Prototype development / 原型开发
  \item Design review / 设计评审
  \item Market test / 市场测试
  \item Product introduction / 产品导入
\end{enumerate}

\term{Concurrent engineering}{并行工程}: design and operations people work together early, avoiding the sequential ``over the wall'' problem.\\
\term{DFM}{可制造性设计}: design products so they can be produced and assembled economically.
\source{Class 7, slides 5-18}

\subsection{质量设计工具 / Quality Design Tools}
\label{subsec:quality-tools}

\begin{longtable}{@{}L{0.18\linewidth}L{0.38\linewidth}L{0.36\linewidth}@{}}
\toprule
\textbf{Tool} & \textbf{中文用途} & \textbf{English use} \\
\midrule
FMEA\index{FMEA} & 系统分析 failure mode, cause, effect, corrective action。 & Systematically analyzes failure modes, causes, effects, and corrective actions. \\
FTA\index{FTA} & 用树形逻辑分析故障之间的关系。 & Visual method for analyzing interrelationships among failures. \\
VA\index{Value analysis} & 问 ``能不能去掉、有没有过度、能不能更便宜更好更快''。 & Eliminates unnecessary features or functions. \\
QFD\index{QFD} & 把顾客之声转化为设计要求。 & Carries the voice of customer into design and implementation. \\
House of Quality\index{House of Quality} & QFD 的第一个矩阵，连接 customer requirements 与 design requirements。 & Matrix connecting customer requirements, technical characteristics, benchmarks, and targets. \\
Kano model\index{Kano model} & 区分 basic, performance, excitement quality。 & Classifies basic, performance, and excitement requirements. \\
\bottomrule
\end{longtable}

\textbf{Kano categories / Kano 三类需求}
\begin{itemize}
  \item Basic: presence has limited satisfaction effect, absence causes dissatisfaction.
  \item Performance: satisfaction increases or decreases with performance level.
  \item Excitement: unexpected feature that delights the customer.
\end{itemize}

\subsection{服务设计 / Service Design}
\label{subsec:service-design}
\index{Service design}
\index{Service blueprinting}

\term{Service concept}{服务概念}: purpose of the service, target market, and desired customer experience.\\
\term{Service package}{服务包}: mixture of physical items, sensual benefits, and psychological benefits.\\
\term{Service specifications}{服务规格}: performance specifications translated into design and delivery specifications.\\
\source{class 8, slides 11-14}

\textbf{Service blueprinting steps / 服务蓝图步骤}
\begin{enumerate}
  \item Establish boundaries.
  \item Identify sequence of customer interactions.
  \item Develop time estimates.
  \item Identify potential failure points.
  \item Establish a time frame.
  \item Analyze profitability.
\end{enumerate}

\begin{longtable}{@{}L{0.30\linewidth}L{0.30\linewidth}L{0.32\linewidth}@{}}
\toprule
\textbf{Design decision} & \textbf{High-contact service} & \textbf{Low-contact service} \\
\midrule
Facility location & Convenient to customers. & Near labor or transportation source. \\
Facility layout & Presentable, supports interaction. & Designed for efficiency. \\
Quality control & More variable; customer participates. & Standards, testing, rework possible. \\
Capacity & Needs excess capacity for peaks. & Planned closer to average demand. \\
Worker skills & Interaction and judgment. & Technical skills. \\
Process & Mostly front-room activities. & Mostly back-room activities. \\
\bottomrule
\end{longtable}

\section{决策理论 / Decision Theory}
\label{sec:decision}

\subsection{决策要素与环境 / Elements and Environments}
\index{Decision theory}
\term{State of nature}{自然状态}: possible future condition that affects payoff.\\
\term{Alternative}{备选方案}: action available to the manager.\\
\term{Payoff}{收益}: outcome of an alternative under a state of nature.\\
\source{Class 9, slides 3-6}

\begin{longtable}{@{}L{0.24\linewidth}L{0.34\linewidth}L{0.34\linewidth}@{}}
\toprule
\textbf{Environment} & \textbf{中文} & \textbf{English} \\
\midrule
Certainty & 结果确定。 & Relevant outcomes are known. \\
Uncertainty & 自然状态概率未知。 & Probabilities are unknown. \\
Risk & 自然状态概率已知或可估计。 & Probabilities are known or estimated. \\
\bottomrule
\end{longtable}

\subsection{不确定性决策 / Decision Under Uncertainty}
\label{subsec:uncertainty}
\index{Maximin}
\index{Maximax}
\index{Laplace}
\index{Minimax regret}

\begin{longtable}{@{}L{0.22\linewidth}L{0.42\linewidth}L{0.28\linewidth}@{}}
\toprule
\textbf{Criterion} & \textbf{Rule / 规则} & \textbf{Mindset} \\
\midrule
Maximin & For each alternative, find worst payoff; choose the best worst payoff. & Conservative / 保守 \\
Maximax & For each alternative, find best payoff; choose the best best payoff. & Optimistic / 乐观 \\
Laplace & Average payoffs assuming equal probabilities; choose highest average. & Equal likelihood / 等可能 \\
Minimax regret & Convert to regret table; for each alternative find maximum regret; choose smallest maximum regret. & Avoid regret / 避免后悔 \\
\bottomrule
\end{longtable}

\exam{题干出现 ``best worst'' 就是 maximin。题干出现 ``least regrets'' 就是 minimax regret。}

\subsection{风险下决策与 EMV / Decision Under Risk and EMV}
\label{subsec:emv}
\index{EMV}
\[
EMV_i=\sum_s p_s \times payoff_{i,s}
\]
\cn{对每个方案分别计算 EMV，选择最大值。}\\
\en{Compute EMV for each alternative and choose the highest value.}\\
\source{Class 9, slides 27-29}

\subsection{决策树 / Decision Trees}
\label{subsec:tree}
\index{Decision tree}
\begin{itemize}
  \item Read the tree from left to right.
  \item Analyze the tree from right to left.
  \item Chance node: calculate expected value.
  \item Decision node: choose the best branch.
\end{itemize}
\exam{开卷时不要从左往右直接算。决策树的计算方向一定是从右往左 rollback。}

\subsection{完全信息期望价值 / EVPI}
\label{subsec:evpi}
\index{EVPI}
\[
A=\sum_s p_s \max_i(payoff_{i,s}), \qquad B=\max_i EMV_i
\]
\[
EVPI=A-B
\]
\cn{$A$ 是完美信息下每个自然状态都选最优方案后的期望值；$B$ 是没有完美信息时的最佳 EMV。}\\
\en{$A$ is the expected value with perfect information; $B$ is the best EMV without perfect information.}

\subsection{敏感性分析 / Sensitivity Analysis}
\index{Sensitivity analysis}
\cn{敏感性分析确定概率或参数在什么范围内变化时，当前最优方案仍然最优。}\\
\en{Sensitivity analysis finds the probability or parameter range over which the current best alternative remains best.}

\section{产能规划 / Capacity Planning}
\label{sec:capacity}

\subsection{产能度量 / Capacity Measures}
\label{subsec:cap-measures}
\index{Capacity planning}
\index{Design capacity}
\index{Effective capacity}

\begin{longtable}{@{}L{0.24\linewidth}L{0.34\linewidth}L{0.34\linewidth}@{}}
\toprule
\textbf{Measure} & \textbf{中文} & \textbf{English} \\
\midrule
Design capacity & 设计最大产出率或服务能力。 & Maximum designed output rate or service capacity. \\
Effective capacity & 设计产能扣除维护、休息、废品等 allowance 后的能力。 & Design capacity minus allowances such as personal time, maintenance, and scrap. \\
Actual output & 实际实现的产出率。 & Output actually achieved. \\
\bottomrule
\end{longtable}

\[
\text{Efficiency}=\frac{\text{Actual Output}}{\text{Effective Capacity}},
\qquad
\text{Utilization}=\frac{\text{Actual Output}}{\text{Design Capacity}}
\]
\[
\text{Capacity Cushion}=100\%-Utilization
\]
\source{Class 10, slides 14-23}

\exam{最容易错的是分母：efficiency 用 effective capacity；utilization 用 design capacity。}

\subsection{产能战略 / Capacity Strategy}
\label{subsec:cap-strategy}
\index{Capacity cushion}
\index{Leading strategy}
\index{Following strategy}
\index{Tracking strategy}

\begin{longtable}{@{}L{0.20\linewidth}L{0.38\linewidth}L{0.34\linewidth}@{}}
\toprule
\textbf{Strategy} & \textbf{中文} & \textbf{Trade-off} \\
\midrule
Leading & 在需求增长前提前建产能。 & 服务水平高，闲置风险高。 \\
Following & 需求超过当前产能后再扩张。 & 闲置少，可能失去顾客。 \\
Tracking & 小幅多次增加产能，贴近需求。 & 折中，管理复杂度较高。 \\
\bottomrule
\end{longtable}

\textbf{Capacity planning steps / 产能规划步骤}
\begin{enumerate}
  \item Estimate future capacity requirements.
  \item Evaluate existing capacity.
  \item Identify alternatives.
  \item Conduct financial analysis.
  \item Assess qualitative issues.
  \item Select one alternative.
  \item Implement.
  \item Monitor results.
\end{enumerate}

\subsection{加工需求 / Processing Requirements}
\index{Processing requirement}
\[
NR=\frac{\sum_{i=1}^{k}p_iD_i}{T}
\]
\begin{itemize}
  \item $p_i$: processing time per unit of product $i$.
  \item $D_i$: demand for product $i$.
  \item $T$: available work time per resource.
  \item If capacity cushion is required, multiply required good output by $1+\text{cushion}$ before calculating resource needs.
\end{itemize}

\subsection{服务产能、瓶颈与约束 / Service Capacity, Bottlenecks, and Constraints}
\label{subsec:bottleneck}
\index{Bottleneck}
\index{Constraint management}

\textbf{Service capacity特点 / Service capacity features}
\begin{itemize}
  \item Need to be near customers; capacity and location are closely tied.
  \item Services cannot be stored; timing must match demand.
  \item Demand is often volatile with peak periods.
\end{itemize}

\term{Bottleneck}{瓶颈}: an operation with the lowest effective capacity that limits total system output.\\
\exam{在 bottleneck 之前增加产能通常不会提高系统产出，因为系统被 bottleneck 限制。}

\textbf{Constraint management / 约束管理}
\begin{enumerate}
  \item Identify the most pressing constraint.
  \item Change operation to maximize benefit from the constraint.
  \item Make sure other parts support the constraint.
  \item Explore and evaluate alternatives.
  \item Repeat when the constraint shifts.
\end{enumerate}

\subsection{成本-产量分析 / Cost-Volume Analysis}
\label{subsec:cvp}
\index{Cost-volume analysis}
\index{Break-even point}

\[
TC=FC+vcQ,\qquad TR=RQ
\]
\[
Profit=TR-TC=Q(R-vc)-FC
\]
\[
Q_{BEP}=\frac{FC}{R-vc}
\]
\[
Q_{\text{target profit}}=\frac{P_{\text{target}}+FC}{R-vc}
\]
\source{Class 11, slides 23-39}

\exam{Make-or-buy 题：分别写 $TC_{make}$ 和 $TC_{buy}$，令二者相等求临界产量。低于临界点通常买，高于临界点通常自制，但要看固定成本在哪个方案。}

\textbf{Assumptions / 假设}
\begin{itemize}
  \item One product is involved.
  \item Everything produced can be sold.
  \item Variable cost per unit is constant.
  \item Fixed cost does not change within the relevant range.
  \item Revenue per unit is constant.
\end{itemize}

\section{流程选择与设施布局 / Process Selection and Facility Layout}
\label{sec:layout}

\subsection{流程类型 / Process Types}
\label{subsec:process-types}
\index{Process selection}
\index{Job shop}
\index{Batch}
\index{Continuous flow}

\begin{longtable}{@{}L{0.22\linewidth}L{0.34\linewidth}L{0.36\linewidth}@{}}
\toprule
\textbf{Type} & \textbf{中文} & \textbf{Key position} \\
\midrule
Job shop & 小规模、高品种、高柔性。 & Low volume, high variety, high equipment flexibility. \\
Batch & 中等产量和中等品种。 & Moderate volume and variety. \\
Repetitive / assembly line & 高产量标准化产品或服务。 & High volume, low variety, standardized flow. \\
Continuous flow & 极高产量、非离散产品。 & Very high volume, very low variety. \\
Project & 非常规、一次性任务。 & One-of-a-kind, often fixed-position. \\
\bottomrule
\end{longtable}
\source{Class 12, slides 11-15}

\subsection{自动化 / Automation}
\index{Automation}
\begin{longtable}{@{}L{0.22\linewidth}L{0.36\linewidth}L{0.34\linewidth}@{}}
\toprule
\textbf{Type} & \textbf{中文} & \textbf{English} \\
\midrule
Fixed automation & 高产量低变化，重大变化成本高。 & Low cost at high volume, minimal variety, costly to change. \\
Programmable automation & 通用设备由计算机程序控制，适合小批量多品种。 & Computer-controlled general equipment for low-volume variety. \\
Flexible automation & 更广品种、中高产量，结合 CNC/CAM/CIM。 & Economically produces wider variety in medium to large volumes. \\
\bottomrule
\end{longtable}

\textbf{Pros and cons / 优缺点}
\begin{itemize}
  \item Pros: low variability, fewer labor issues, lower variable cost.
  \item Cons: high fixed cost, barrier to change, possible adverse morale/productivity effects.
\end{itemize}

\subsection{基础布局类型 / Basic Layout Types}
\label{subsec:layouts}
\index{Product layout}
\index{Process layout}
\index{Fixed-position layout}
\index{Cellular layout}

\begin{longtable}{@{}L{0.22\linewidth}L{0.34\linewidth}L{0.36\linewidth}@{}}
\toprule
\textbf{Layout} & \textbf{中文} & \textbf{English exam cue} \\
\midrule
Product layout & 按产品加工顺序排列，适合标准化、高产量。 & Sequential arrangement, high output, low unit cost, vulnerable to shutdowns. \\
Process layout & 按功能/部门分组，适合多样化加工。 & Functional grouping, flexible, routing/scheduling/material handling complex. \\
Fixed-position layout & 产品/项目不移动，人、材料、设备移动。 & Product stays stationary due to weight, size, or bulk. \\
Cellular layout & 把机器组成 cell，处理相似零件族。 & Group technology, combines flow improvement with some flexibility. \\
\bottomrule
\end{longtable}
\source{Class 12, slides 28-39; Class 13, slides 4-9}

\exam{若题目说 hospital，一般偏 process layout；若说 automobile assembly，偏 product layout；若说 infrastructure project，偏 fixed-position layout。}

\subsection{生产线平衡 / Line Balancing}
\label{subsec:line-balancing}
\index{Line balancing}
\index{Cycle time}

\term{Line balancing}{生产线平衡}: assigning tasks to workstations so that workstation time requirements are approximately equal.\\
\source{Class 13, slides 11-30}

\[
\text{Output rate}=\frac{OT}{CT}, \qquad CT=\frac{OT}{D}
\]
\[
N_{\min}=\left\lceil\frac{\sum t_i}{CT}\right\rceil
\]
\[
\text{Percent idle time}=\frac{\text{Idle time per cycle}}{N\times CT}, \qquad
\text{Efficiency}=1-\text{Percent idle time}
\]

\textbf{Assignment procedure / 分配步骤}
\begin{enumerate}
  \item Determine cycle time and minimum number of workstations.
  \item Draw or read the precedence diagram.
  \item Start with Station 1 and move left to right.
  \item Eligible task must have all predecessors assigned and must fit in remaining station time.
  \item If no eligible task fits, move to next workstation.
  \item After all tasks are assigned, compute idle time and efficiency.
\end{enumerate}

\textbf{Heuristic rules / 启发式规则}
\begin{itemize}
  \item Assign tasks in order of most following tasks.
  \item Assign tasks in order of greatest positional weight.
  \item Tie-breaker: longest task time or more followers.
\end{itemize}

\subsection{流程布局运输成本 / Process Layout Load-Distance Cost}
\label{subsec:load-distance}
\index{Load-distance cost}
\[
\text{Total cost}=\sum_{i<j}F_{ij}D_{ij}C_{ij}
\]
\begin{itemize}
  \item $F_{ij}$: work flow between departments $i$ and $j$.
  \item $D_{ij}$: distance between assigned locations.
  \item $C_{ij}$: unit movement cost. If not given, treat as 1.
\end{itemize}
\exam{核心直觉：物流量大的部门放近，物流量小的部门可以放远。}

\section{题型模板 / Exam Calculation Templates}
\label{sec:templates}

\subsection{预测计算模板 / Forecasting Template}
\begin{enumerate}
  \item Identify the method: naive, MA, WMA, exponential smoothing, trend, or regression.
  \item Compute forecast values.
  \item Build error table: $A$, $F$, $A-F$, $|A-F|$, $(A-F)^2$, $\left|\frac{A-F}{A}\right|100$.
  \item Compute MAD, MSE, MAPE.
  \item Select the method with lower error, unless the question asks for a specific method.
\end{enumerate}

\subsection{决策表模板 / Decision Table Template}
\begin{enumerate}
  \item Draw payoff table: alternatives as rows, states of nature as columns.
  \item If probabilities are unknown, choose maximin/maximax/Laplace/minimax regret.
  \item If probabilities are known, compute EMV for every alternative.
  \item For EVPI, compute $A$ under perfect information and $B$ as best EMV.
\end{enumerate}

\subsection{产能与 BEP 模板 / Capacity and Break-even Template}
\begin{enumerate}
  \item Mark design capacity, effective capacity, actual output.
  \item Compute efficiency and utilization with correct denominator.
  \item For capacity cushion, compute $100\%-utilization$.
  \item For CVP, write $TC=FC+vcQ$ and $TR=RQ$.
  \item Solve $TR=TC$ for BEP, or $Profit=P_{target}$ for target-profit volume.
\end{enumerate}

\subsection{生产线平衡模板 / Line Balancing Template}
\begin{enumerate}
  \item Compute $CT=OT/D$.
  \item Compute $N_{\min}=\lceil \sum t_i/CT\rceil$.
  \item Draw precedence diagram.
  \item Assign tasks by eligibility and heuristic rule.
  \item Compute total available station time $N\times CT$.
  \item Compute idle time, percent idle time, and efficiency.
\end{enumerate}

\section{高频易错点 / Common Pitfalls}
\label{sec:pitfalls}

\begin{enumerate}
  \item \kw{Efficiency vs. utilization}: efficiency denominator is effective capacity; utilization denominator is design capacity.
  \item \kw{MAD vs. MSE}: MAD is not square root of MSE.
  \item \kw{Tracking signal}: use cumulative forecast error divided by MAD; large absolute value signals bias.
  \item \kw{Maximin}: best of worst, not worst of best.
  \item \kw{Minimax regret}: first build regret table by subtracting each payoff from the best payoff in its state column.
  \item \kw{Decision tree}: read left to right, calculate right to left.
  \item \kw{BEP}: denominator is contribution margin $R-vc$, not $R/VC$.
  \item \kw{Capacity cushion}: if utilization is 72\%, cushion is 28\%.
  \item \kw{Line balancing}: tasks cannot be assigned before predecessors are assigned, even if they fit time.
  \item \kw{Process vs. product layout}: process layout is functional grouping; product layout is sequential flow.
  \item \kw{Delayed differentiation}: means postponement of final customization, not full customization from the beginning.
  \item \kw{Kano}: basic features prevent dissatisfaction but do not create much satisfaction when present.
\end{enumerate}

\section{一页术语表 / One-page Glossary}
\label{sec:glossary}

\begin{multicols}{2}
\begin{itemize}
  \item Operations management - 运营管理
  \item Transformation process - 转化过程
  \item System design decisions - 系统设计决策
  \item System operation decisions - 系统运行决策
  \item Trade-off - 权衡
  \item Competitiveness - 竞争力
  \item Order qualifiers - 订单资格要素
  \item Order winners - 赢单要素
  \item Productivity - 生产率
  \item Forecasting - 预测
  \item Moving average - 移动平均
  \item Exponential smoothing - 指数平滑
  \item MAD - 平均绝对偏差
  \item MSE - 均方误差
  \item MAPE - 平均绝对百分比误差
  \item Tracking signal - 跟踪信号
  \item Standardization - 标准化
  \item Delayed differentiation - 延迟差异化
  \item Modular design - 模块化设计
  \item Reliability - 可靠性
  \item Robust design - 稳健设计
  \item FMEA - 失效模式与影响分析
  \item FTA - 故障树分析
  \item QFD - 质量功能展开
  \item House of Quality - 质量屋
  \item Kano model - Kano 模型
  \item Decision theory - 决策理论
  \item State of nature - 自然状态
  \item Maximin - 最大最小准则
  \item Maximax - 最大最大准则
  \item Minimax regret - 最小最大后悔准则
  \item EMV - 期望货币价值
  \item EVPI - 完全信息期望价值
  \item Design capacity - 设计产能
  \item Effective capacity - 有效产能
  \item Utilization - 利用率
  \item Efficiency - 效率
  \item Capacity cushion - 产能缓冲
  \item Bottleneck - 瓶颈
  \item Break-even point - 盈亏平衡点
  \item Job shop - 作业车间
  \item Batch - 批量流程
  \item Continuous flow - 连续流程
  \item Product layout - 产品布局
  \item Process layout - 流程布局
  \item Fixed-position layout - 固定位置布局
  \item Cellular layout - 单元布局
  \item Line balancing - 生产线平衡
\end{itemize}
\end{multicols}

\clearpage
\renewcommand{\indexname}{LaTeX 术语索引 / LaTeX Term Index}
\addcontentsline{toc}{section}{LaTeX 术语索引 / LaTeX Term Index}
\printindex

\end{document}
